% !TeX spellcheck = en_US
\documentclass[french]{yLectureNote}

\title{Électromagnétisme 2 PS}
\subtitle{Physique}
\author{Paulhenry Saux}
\date{\today}
\yLanguage{Français}
%lionels.calmels@cemes.fr
\professor{M.Brut}
\usepackage{graphicx}%----pour mettre des images
\usepackage[utf8]{inputenc}%---encodage
\usepackage{geometry}%---pour modifier les tailles et mettre a4paper
%\usepackage{awesomebox}%---pour les boites d'exercices, de pbq et de croquis ---d\'esactiv\'e pour les TP de PC
\usepackage{tikz}%---pour deiffner + d\'ependance de chemfig
% \usepackage{tabularx}%---pour dimensionner automatiquement les tableaux avec variable X
\usepackage{awesomebox}%---Pour les boites info, danger et autres
\usepackage{menukeys}%---Pour deiffner les touches de Calculatrice
\usepackage{fancyhdr}%---pour les en-t\^ete personnalis\'ees
\usepackage{blindtext}%---pour les liens
\usepackage{hyperref}%---pour les liens (\`a mettre en dernier)
\usepackage{caption}%---pour la francisation de la l\'egende table vers Tableau
\usepackage{pifont}
\usepackage{array}%---pour les tableaux
\usepackage{lipsum}
\usepackage{yFlatTable}
\usepackage{multicol}
\usepackage{tabularx}
\usepackage{booktabs}
\usepackage{esint}% Pour les oint
\newcommand{\Lim}[1]{\lim\limits_{\substack{#1}}\:}
\renewcommand{\vec}{\overrightarrow}
\newcommand{\N}[0]{\mathbb{N}}
\newcommand{\dd}{\mathrm{d}}
\newcommand{\norm}[1]{||\vec{#1}||}
\newcommand{\fo}{\psi(\vec{r},t)}
\newcommand{\foe}{\psi(\vec{r},t)\*}
\newcommand{\HH}{\hat{H}}
\newcommand{\hb}{\hbar}
\newcommand{\lap}{\nabla^2}
\newcommand{\lapcc}{\frac{\partial^2 }{\partial x^2}+\frac{\partial^2 }{\partial y^2}+\frac{\partial^2 }{\partial z^2}}
\newcommand{\E}{\vec{E}(M)}
\newcommand{\B}{\vec{B}(M)}
\newcommand{\V}{V(M)}
\newcommand{\er}{\vec{e_{\rho}}}
\newcommand{\ep}{\vec{e_{\varphi}}}
\newcommand{\pip}{\pi^+}
\newcommand{\pim}{\pi^-}
\newcommand{\mv}{\mathcal{V}}
\DeclareMathOperator{\Div}{Div}
\newcommand{\rot}{\vec{\mathrm{rot}}}
\newcommand{\grad}{\vec{\mathrm{grad}}}
\setcounter{chapter}{3}
\begin{document}
\chapter{Puissance et énergie du champ Électro-magnétique}

\section{Puissance et Théorème de Poynting}

\begin{definition}[Puissance volumique de Lorentz]
La puissance transférée par le champ EM à la matière (densité de courant $\vec{j}$) est :
$$P_{vol} = \vec{j} \cdot \vec{E}$$
\end{definition}

\tipsInfo{Sens Physique}{La force magnétique ne travaille jamais. Le transfert d'énergie ne se fait que via le champ électrique. Si $\vec{j} \cdot \vec{E} > 0$, le champ cède de l'énergie (effet Joule, accélération de charges). Si $\vec{j} \cdot \vec{E} < 0$, la matière fournit de l'énergie au champ (générateur).}

\begin{theorem}[Théorème de Poynting (Local)]
L'énergie électromagnétique obéit à une équation de conservation locale :
$$-\vec{j} \cdot \vec{E} = \frac{\partial u_{em}}{\partial t} + \Div \vec{\Pi}$$
\end{theorem}
\begin{definition}[Densité volumique d'énergie EM]
    $u_{em} = \frac{\epsilon_0 E^2}{2} + \frac{B^2}{2\mu_0}$ (en $J.m^{-3}$)
\end{definition}
\begin{definition}[Vecteurs de Poyting]
    $\vec{\Pi} = \frac{\vec{E} \wedge \vec{B}}{\mu_0}$ (en $W.m^{-2}$)
\end{definition}

Le terme $\Div \vec{\Pi}$ représente la puissance sortante par unité de volume. Le vecteur $\vec{\Pi}$ indique la direction et l'intensité du flux d'énergie électromagnétique.

\section{Bilan Énergétique Intégral}

\begin{proposition}[Bilan sur un volume $V$]
$$\frac{\dd U_{em}}{\dd t} = - \oiint_{S} \vec{\Pi} \cdot \dd \vec{S} - P_{meca}$$
\end{proposition}



\warningInfo{Signe du flux}{Par convention, $\dd \vec{S}$ est sortant. Donc $\oiint \vec{\Pi} \cdot \dd \vec{S} > 0$ signifie que de l'énergie quitte le volume (rayonnement).}

\section{Applications aux Composants}

\begin{table}[h]
\centering
\begin{tabularx}{\textwidth}{|l|X|X|}
\hline
\textbf{Composant} & \textbf{Énergie Stockée} & \textbf{Origine du flux} \\ \hline
\textbf{Condensateur} & $U_e = \frac{1}{2}CU^2$ & Énergie stockée dans $\vec{E}$ entre les armatures. \\ \hline
\textbf{Inductance} & $U_m = \frac{1}{2}LI^2$ & Énergie stockée dans $\vec{B}$ au cœur des spires. \\ \hline
\textbf{Résistance} & $U_{em} \approx 0$ (statique) & Énergie dissipée ($RI^2$) venant de l'extérieur du fil. \\ \hline
\end{tabularx}
\end{table}

% \subsection{Le "Paradoxe" de la Résistance}

% \criticalInfo{Le flux entrant}{Dans une résistance cylindrique, $\vec{E}$ est axial et $\vec{B}$ est orthoradial à la surface. Le vecteur de Poynting $\vec{\Pi} = \frac{\vec{E} \wedge \vec{B}}{\mu_0}$ pointe \textbf{vers l'axe du fil}.}



% \begin{flalign*}
%     &\text{À la surface } (\rho = a) : && \vec{E} = \frac{V}{l}\vec{e}_z \quad \text{et} \quad \vec{B} = \frac{\mu_0 I}{2\pi a}\vec{e}_{\theta} &\\
%     &\text{Flux de Poynting :} && \Phi_{entrant} = -\oiint \vec{\Pi} \cdot \dd \vec{S} = \frac{V}{l} \frac{I}{2\pi a} (2\pi a l) = VI = RI^2 &
% \end{flalign*}

% \tipsInfo{Compréhension}{L'énergie ne circule pas \textit{dans} le fil via les électrons, mais elle est transmise par le champ EM \textit{autour} du fil et "s'engouffre" latéralement dans le conducteur pour y être dissipée.}

\end{document}
